\documentclass{article}

\usepackage{microtype}
\usepackage{graphicx}
\usepackage{subcaption}
\usepackage{booktabs} % for professional tables
\usepackage{hyperref}
\newcommand{\theHalgorithm}{\arabic{algorithm}}
% Course style (ethdl2026, from Moodle); [accepted] shows the author names.
\usepackage[accepted]{ethdl2026}
\usepackage{amsmath}
\usepackage{amssymb}
\usepackage{mathtools}
\usepackage{amsthm}
\usepackage[capitalize,noabbrev]{cleveref}
\theoremstyle{plain}
\newtheorem{theorem}{Theorem}[section]
\newtheorem{proposition}[theorem]{Proposition}
\newtheorem{lemma}[theorem]{Lemma}
\newtheorem{corollary}[theorem]{Corollary}
\theoremstyle{definition}
\newtheorem{definition}[theorem]{Definition}
\newtheorem{assumption}[theorem]{Assumption}
\theoremstyle{remark}
\newtheorem{remark}[theorem]{Remark}
\usepackage[textsize=tiny]{todonotes}
\icmltitlerunning{ETHZ Deep Learning 2026 -- Project 14}

\begin{document}

\twocolumn[
  \icmltitle{ETHZ Deep Learning 2026 Project Proposal \\
    Project 14: How Much of the Learned Testing Policy Should a Model Keep?}

  \icmlsetsymbol{equal}{*}

  % TODO: full names and ETH usernames
  \begin{icmlauthorlist}
    \icmlauthor{Daniele Bugatti}{ETH}
    \icmlauthor{Rafael van 't Veer}{ETH}
    \icmlauthor{Ejtehadi Mehdi}{ETH}
    \icmlauthor{Wesley Remy}{ETH}
  \end{icmlauthorlist}

  \icmlaffiliation{ETH}{Department of Computer Science, ETH Zurich, Zurich, Switzerland}

  % TODO: real e-mail; the course style requires at least one corresponding author
  \icmlcorrespondingauthor{Wesley Remy}{username@ethz.ch}

  \icmlkeywords{Observation shift, informative missingness, clinical time series, state-space models, domain adaptation}

  \vskip 0.3in
]

\printAffiliationsAndNotice{Advisor TA: Sarah Martinson. Topic: Project 14, Deep Time-Series Models under Observation Shift. }

\begin{abstract}
Clinical time-series models learn from \emph{when} and \emph{whether} tests are ordered, not only from their values. This testing pattern reflects one hospital's observation policy and changes when the model is deployed elsewhere. Before deployment, hospitals typically run a model silently for 1--3 months, which yields only a handful of labelled outcome cases. We propose a ladder of adaptations that acts only on the testing part of a latent state-space model, from a single trust weight through per-test-group weights to re-learning the testing head, and ask which rung a run-in of a given size supports for which kind of observation shift.
\end{abstract}

\section{Introduction \& Related Work}

\textbf{Problem.} Deep models for clinical time series learn from what is measured and also from when it is measured: a lactate order signals that a doctor suspects sepsis. This testing pattern is predictive \citep{sisk2021}, but it reflects the hospital's observation policy, and it changes when a model trained at hospital A is deployed at hospital B (\emph{observation shift}). The change comes in two kinds. Hospital B may test everyone more, so every order means a bit less; or individual tests may change meaning: troponin is ordered for 3\% of patients at Beth Israel but 32\% at Emory (PhysioNet 2019, \citealp{reyna2020}), while a lactate order may keep its meaning. A model adapted to B must fix the first kind globally and the second per test.

Before going live, models are often run silently for 1--3 months \citep{sendak2020}. That run-in yields the only labelled data from B, and very little of it: a typical ICU admits 600--1{,}300 patients a year \citep{checkley2014}, so a 1--3 month run-in contains about 2--6 sepsis cases. We therefore propose a \textbf{ladder of adaptations} that changes only the testing part of a latent state-space model in which the testing evidence enters through its own head: (1) one global weight $w$ on the testing evidence ($w{=}0$ ignore, $w{=}1$ trust as at A); (2) one weight per group of tests; (3) re-learning the testing head with the rest frozen; (4) fine-tuning the whole model. Higher rungs repair more kinds of shift but need more outcome cases.

\textbf{Research question.} \emph{For a run-in of a given size at a new hospital, which rung of adaptation of the testing part predicts best, and how does the answer depend on the kind of observation shift?} Concretely: (i) does adapting only the testing part beat the fixed rules of prior work (always keep or always remove the pattern) and full fine-tuning at realistic run-in sizes; (ii) at how many outcome events does each higher rung start to pay off; (iii) does the state-space model adapt better from few cases than a standard architecture (GRU-D) fine-tuned at its last layer or in full.

\textbf{Deep models use the observation pattern but do not separate it.} Architectures for irregular series consume masks and time gaps (GRU-D, \citealp{che2018}), continuous-time paths (Neural CDE, \citealp{kidger2020}; Latent ODE, \citealp{rubanova2019}) or time-value sets with attention (mTAN, \citealp{shukla2021}; Raindrop, \citealp{zhang2022}), and fuse the pattern with physiology in one representation. Models with a separate observation component exist, a Hawkes process for clinicians' measurement events \citep{alaa2017} and DeepJoint's observation heads \citep{jeanselme2022}, but none adapts that component alone at a new site.

\textbf{Observation shift between hospitals.} Multi-site deep sepsis and ICU benchmarks lose 0.08--0.09 AUROC across sites \citep{moor2023,yaib2024} and include missingness indicators without attributing the loss to them. Studies that isolate the pattern choose all or nothing: \citet{futoma2021} fitted with measurement indicators and dropped them at deployment, which transferred best from Beth Israel to a Memphis hospital but not in reverse; \citet{rockenschaub2024} found that ignoring informative missingness costs at the source but is more robust elsewhere, especially when missingness is highly informative. \citet{lasko2024} argue conceptually for a site-specific process model next to a stable disease model; \citet{zhou2023} adapt to missingness shift without labels in linear models. None tunes \emph{how much} of the learned policy to keep from labelled target data.

\textbf{Adapting with few labels.} Clinical model updating is a ladder from recalibration to full refitting, and parsimonious updates win at small samples \citep{steyerberg2004,vergouwe2017}; a full refit is risky well below the 100 events needed just to validate a model \citep{riley2024}. In deep learning, which part to fine-tune depends on the kind of shift \citep{lee2023}, and retraining only the last layer can suffice \citep{kirichenko2023}, but clinical fine-tuning studies use hundreds to thousands of target patients \citep{moor2023,yaib2024}. \textbf{Gap:} no work adapts only the observation part of a clinical sequence model, and none compares such adaptations against run-in size.

\section{Experimental Setup \& Methods}

\textbf{Model.} A neural latent state-space model, extending our pilot: a discrete latent patient state ($K{\approx}32$ states) evolves by a learned transition, and three heads read it: a value head (likelihood of each measured value), a testing head (probability that each test is ordered, given the state) and an outcome head. Predictions come from filtering the state over time, trained end to end on the outcome loss plus the likelihood of values and tests. The testing pattern enters only through the testing head's log-likelihood, which is multiplied by the weight(s) of rungs 1--2; rung 3 re-learns the testing head with transition, value and outcome heads frozen; rung 4 fine-tunes everything. Baselines: keep ($w{=}1$), remove ($w{=}0$, Futoma), recalibration, and a standard GRU-D (fine-tuned at the last layer or in full). Lessons from our pilots are built in: tuned weights are shrunk towards 1, calibration is fitted separately from the weights, and impossible values are removed before recalibration.

\textbf{E1: simulation with known optimum.} Our latent-state simulator, whose exact Bayes predictor is known, with two shift types: global (B tests less selectively) and test-specific (one test group becomes routine). We measure the regret of each rung against the Bayes predictor at B for run-ins of 5--200 outcome events. Expected: rung 1 wins under global shift and small run-ins, rungs 2--3 under test-specific shift once enough events exist.

\textbf{E2: real hospitals (dataset to be determined).} Requirements: several hospitals with hospital identifiers, time-stamped labs and vitals with the testing pattern preserved, and enough outcome events per site; the main candidate is eICU-CRD (208 US hospitals, \citealp{pollard2018}; credentialed access). Design: ordered source $\to$ target pairs, run-ins of 1, 3, 6, 12 and 24 months drawn from the target site (20 draws each), evaluation on all remaining target patients, primary metric AUROC (pre-specified), secondary calibration.

\textbf{E3: architecture.} Whether the state-space model adapts from fewer cases than GRU-D fine-tuned at the last layer or in full, on the pairs of E2.

\textbf{Computational cost and feasibility.} We have one NVIDIA GTX 1080 Ti (11\,GB) on snowflake and a 4-core laptop. All models are small (32 latent states or GRU hidden size 32--64, at most 48 hourly steps per patient), so GPU memory is not a constraint. Measured on the laptop CPU: the pilot state-space model trains in 15 minutes on 8{,}000 simulated patients, and 108 simulation fits of GRU, GRU-D, Neural CDE and Transformer took 3.6 hours in total; E1 therefore runs on the laptop. E2 and E3 need about 50 source models (sites $\times$ models $\times$ 3 seeds) and a few thousand short adaptation runs, which we run on the GPU. We will time one site first and, if the full design does not fit the available weeks, reduce the number of draws or pairs.

\section{Results \& Discussion}

\textbf{Preliminary results} (linear models, open PhysioNet 2019 and 2012 data). On Beth Israel $\to$ Emory, keeping the pattern was best, because there was barely any observation shift between the two (pattern-model transfer gap 0.03 AUROC, 95\% CI $-0.01$ to $0.07$); dropping it cost 0.05 AUROC. Across 26 unit pairs the pattern's meaning shifted clearly in 7, mostly for mortality; there, a single weight tuned on a 3--24 month run-in gained up to 0.02--0.03 AUROC and stayed within 0.001 elsewhere (after removing impossible lab values). A one-parameter update beat a full refit up to about 20 cases; the refit won from about 50. In simulation, the state-space model trained at A and tempered by $w$ followed the exact optimal trade-off between A and B to within 0.001 Brier; it has not yet been run on real data.

\section{Conclusion}

The project turns the keep-or-remove question of prior work into a choice of how much of the testing part to adapt, given the cases a run-in actually yields. The simulator, the pilot model and the run-in pipeline already exist. Timeline: weeks 1--2 the multivariate state-space model and E1; weeks 3--4 E2 and E3; weeks 5--6 analysis and report.

\bibliography{bibliography}
\bibliographystyle{ethdl2026}

%%%%%%%%%%%%%%%%%%%%%%%%%%%%%%%%%%%%%%%%%%%%%%%%%%%%%%%%%%%%%%%%%%%%%%%%%%%%%%%
%%%%%%%%%%%%%%%%%%%%%%%%%%%%%%%%%%%%%%%%%%%%%%%%%%%%%%%%%%%%%%%%%%%%%%%%%%%%%%%
% APPENDIX
%%%%%%%%%%%%%%%%%%%%%%%%%%%%%%%%%%%%%%%%%%%%%%%%%%%%%%%%%%%%%%%%%%%%%%%%%%%%%%%
%%%%%%%%%%%%%%%%%%%%%%%%%%%%%%%%%%%%%%%%%%%%%%%%%%%%%%%%%%%%%%%%%%%%%%%%%%%%%%%
% \newpage
% \appendix
% \onecolumn
% \section{Appendix}

\end{document}
